\begin{lemma}[移位 Fibonacci 的严格融合恒等式与正缺陷项]\label{lem:pom-shifted-fib-fusion-defect-positive}
取 Fibonacci 数列 \(F_0=0,F_1=1\)。
则对任意整数 \(a,b\ge 0\) 恒有恒等式
\begin{equation}\label{eq:pom-shifted-fib-fusion-defect}
\boxed{\;
F_{a+2}F_{b+2}=F_{a+b+2}+F_aF_b.\;}
\end{equation}
特别地，
\(
F_{a+2}F_{b+2}\ge F_{a+b+2}
\)，且严格性当且仅当 \(a,b\ge 1\)；
因此移位序列 \(G(n):=F_{n+2}\) 在加法半群 \((\NN_{\ge 1},+)\) 上严格超乘法，
其“融合缺陷”由正项 \(F_aF_b\) 精确给出。
\end{lemma}
\begin{proof}
由 Fibonacci 加法公式
\(
F_{m+n}=F_mF_{n+1}+F_{m-1}F_n
\)，取 \(m=a+2\)、\(n=b\) 得
\[
F_{a+b+2}=F_{a+2}F_{b+1}+F_{a+1}F_b.
\]
又由递推 \(F_{b+2}=F_{b+1}+F_b\) 与 \(F_{a+2}=F_{a+1}+F_a\)，有
\[
F_{a+2}F_{b+2}
=F_{a+2}(F_{b+1}+F_b)
=\underbrace{F_{a+2}F_{b+1}+F_{a+1}F_b}_{=\,F_{a+b+2}}
\;+\;F_aF_b,
\]
即得 \eqref{eq:pom-shifted-fib-fusion-defect}。
当且仅当 \(F_aF_b>0\)（即 \(a,b\ge 1\)）时为严格不等式。证毕。
\end{proof}

\begin{corollary}[二维行列式的 Fibonacci 体积律]\label{cor:pom-fib-determinant-volume-law}
对任意整数 \(a,b\ge 0\) 定义矩阵
\[
M(a,b):=\begin{pmatrix}
F_{a+2} & F_a\\
F_b & F_{b+2}
\end{pmatrix}.
\]
则其行列式满足严格恒等式
\[
\boxed{\;\det M(a,b)=F_{a+b+2}.\;}
\]
该恒等式把“移位乘积—融合缺陷”结构提升为一个确定性的面积（体积）律：二分量参数对 \((a,b)\) 在整数格上的面积像由单一 Fibonacci 数 \(F_{a+b+2}\) 完全刻画。
\end{corollary}
\begin{proof}
由引理 \ref{lem:pom-shifted-fib-fusion-defect-positive}，
\(
F_{a+2}F_{b+2}-F_aF_b=F_{a+b+2}
\)，而左端即 \(\det M(a,b)\)。证毕。
\end{proof}

\begin{proposition}[路径分量乘法多重度的细化单调性与极值界]\label{prop:pom-path-component-multiplicity-refinement-monotone-extrema}
定义路径分量长度 \(\ell_1,\dots,\ell_r\ge 1\) 的乘法多重度
\[
d(\ell_1,\dots,\ell_r):=\prod_{i=1}^r F_{\ell_i+2},
\qquad
L:=\sum_{i=1}^r \ell_i.
\]
若将两分量 \((a,b)\) 融合为单分量 \((a+b)\)（其中 \(a,b\ge 1\)），则多重度严格下降且下降量具有闭式：
\[
\boxed{\;
d(\ldots,a,b,\ldots)-d(\ldots,a+b,\ldots)
=\Big(\prod_{\text{其余 }i}F_{\ell_i+2}\Big)\,F_aF_b\ >\ 0.\;}
\]
因此 \(d\) 对“分量细化偏序”严格单调：任何细化操作（把某分量拆成两段正长度之和）都严格增大 \(d\)。
从而对任意分解 \(L=\sum_i\ell_i\) 恒有极值界
\[
\boxed{\;
F_{L+2}\ \le\ \prod_{i=1}^r F_{\ell_i+2}\ \le\ 2^L.\;}
\]
左端等号当且仅当 \(r=1\)（单分量）；右端等号当且仅当 \(\ell_i\equiv 1\)（全为长度 \(1\) 的最细分量）。
\end{proposition}
\begin{proof}
差值闭式由引理 \ref{lem:pom-shifted-fib-fusion-defect-positive} 取 \((a,b)\) 并乘以其余因子得到：
\[
F_{a+2}F_{b+2}-F_{a+b+2}=F_aF_b>0\qquad(a,b\ge 1).
\]
\par
\textbf{下界：}从任意分解 \((\ell_1,\dots,\ell_r)\) 出发，逐步融合两分量直至得到单分量 \((L)\)；
每一步融合都使 \(d\) 严格下降，故
\(
d(\ell_1,\dots,\ell_r)\ge d(L)=F_{L+2}
\)，且仅当 \(r=1\) 时取等号。
\par
\textbf{上界：}对每个 \(\ell\ge 1\) 用归纳可证 \(F_{\ell+2}\le 2^\ell\)：\(\ell=1\) 时 \(F_3=2\)；
若对 \(\ell-1,\ell-2\) 成立，则
\(
F_{\ell+2}=F_{\ell+1}+F_\ell\le 2^{\ell-1}+2^{\ell-2}\le 2^\ell
\)。
因此
\(
\prod_i F_{\ell_i+2}\le \prod_i 2^{\ell_i}=2^{\sum_i\ell_i}=2^L
\)。
若某 \(\ell_i\ge 2\)，则 \(F_{\ell_i+2}<2^{\ell_i}\)，从而整体严格小于 \(2^L\)；
故上界取等号当且仅当 \(\ell_i=1\) 对所有 \(i\) 成立。证毕。
\end{proof}

\begin{proposition}[多重度的二参数 Binet 正规形与全局误差界]\label{prop:pom-multiplicity-binet-two-parameter-normal-form}
令 \(\varphi=\frac{1+\sqrt5}{2}\)，并设 \(\ell_1,\dots,\ell_r\ge 1\)、\(L=\sum_i\ell_i\)。
则乘法多重度 \(d(\ell_1,\dots,\ell_r)=\prod_i F_{\ell_i+2}\) 具有精确分解
\[
\boxed{\;
d(\ell_1,\dots,\ell_r)
=\Bigl(\frac{\varphi^2}{\sqrt5}\Bigr)^r\varphi^{L}\prod_{i=1}^r\bigl(1-(-\varphi^{-2})^{\ell_i+2}\bigr).\;}
\]
由于 \(\ell_i\ge 1\Rightarrow |(-\varphi^{-2})^{\ell_i+2}|\le \varphi^{-6}\)，从而得到统一双边界
\[
\boxed{\;
\Bigl(\frac{\varphi^2}{\sqrt5}\Bigr)^r\varphi^{L}(1-\varphi^{-6})^r
\ \le\
d(\ell_1,\dots,\ell_r)
\ \le\
\Bigl(\frac{\varphi^2}{\sqrt5}\Bigr)^r\varphi^{L}(1+\varphi^{-6})^r.\;}
\]
因此 \(d\) 在对数尺度上由两参数 \((L,r)\) 主导，而振荡项由显式指数 \(\varphi^{-6}\) 一致控制。
\end{proposition}
\begin{proof}
由 Binet 公式 \(F_n=(\varphi^n-\psi^n)/\sqrt5\) 且 \(\psi=-\varphi^{-1}\)，得
\[
F_{\ell+2}=\frac{\varphi^{\ell+2}}{\sqrt5}\Bigl(1-\Bigl(\frac{\psi}{\varphi}\Bigr)^{\ell+2}\Bigr)
=\frac{\varphi^{\ell+2}}{\sqrt5}\bigl(1-(-\varphi^{-2})^{\ell+2}\bigr).
\]
对 \(i\) 取乘积得到所示精确分解。
又因 \(\ell_i+2\ge 3\)，有
\(
|(-\varphi^{-2})^{\ell_i+2}|\le \varphi^{-6}
\)，
故每一因子落在区间 \([1-\varphi^{-6},\,1+\varphi^{-6}]\) 内；再取乘积即得双边界。证毕。
\end{proof}

\begin{proposition}[融合缺陷的结合律约束：Fibonacci \texorpdfstring{$2$}{2}-余因子恒等式]\label{prop:pom-fusion-defect-2cocycle-identity}
令 \(G(n):=F_{n+2}\)，并定义融合缺陷项
\(
\omega(a,b):=F_aF_b
\)（\(a,b\ge 0\)）。
由引理 \ref{lem:pom-shifted-fib-fusion-defect-positive} 的融合律
\(
G(a)G(b)=G(a+b)+\omega(a,b)
\)
与乘法结合律，推出对任意 \(a,b,c\ge 0\) 必有严格恒等式
\[
\boxed{\;
\omega(a,b)G(c)+\omega(a+b,c)
=\omega(b,c)G(a)+\omega(b+c,a)
=\omega(c,a)G(b)+\omega(c+a,b).\;}
\]
该恒等式表明 \(\omega\) 作为“融合缺陷项”满足非平凡的 \(2\)-余因子一致性：任意三分量的总缺陷与融合顺序无关，并由上述三种显式表达共同刻画。
\end{proposition}
\begin{proof}
对乘积 \(G(a)G(b)G(c)\) 分别按 \((ab)c\) 与 \(a(bc)\) 两种括号方式展开：
\[
\bigl(G(a)G(b)\bigr)G(c)=\bigl(G(a+b)+\omega(a,b)\bigr)G(c)=G(a+b+c)+\omega(a+b,c)+\omega(a,b)G(c),
\]
\[
G(a)\bigl(G(b)G(c)\bigr)=G(a)\bigl(G(b+c)+\omega(b,c)\bigr)=G(a+b+c)+\omega(a,b+c)+\omega(b,c)G(a),
\]
比较两式并消去公共项 \(G(a+b+c)\) 得
\[
\omega(a,b)G(c)+\omega(a+b,c)
=\omega(b,c)G(a)+\omega(a,b+c).
\]
由于 \(\omega\) 对称（\(\omega(a,b)=\omega(b,a)\)），有 \(\omega(a,b+c)=\omega(b+c,a)\)，从而得到第一条等式。
循环置换 \((a,b,c)\) 得三者相等。证毕。
\end{proof}

\endinput
